\chapter{Topologi Umum}\label{ch:b3:topology}

Jilid Tahun ke-2 melakukan analisis di ruang metrik: jarak,
bola, barisan. Namun gagasan yang mendasar --- \omterm{def:b3:topology:continuity}{kekontinuan},
kekompakan, keterhubungan --- tak pernah menyebut nilai numerik
sebuah jarak, melainkan hanya keluarga \emph{\omterm{def:b3:topology:topology}{himpunan terbuka}} yang dibangkitkannya.
Bab ini mengambil keluarga itu sebagai objek primitifnya. Untungnya
bukan keumuman demi keumuman: konstruksi kuosien (misalnya
lingkaran sebagai $\R/\Z$, ruang proyektif), hasil kali, dan
\omterm{def:b3:topology:topology}{topologi} bertipe lemah pada analisis fungsional memang bukan
objek yang metrik-dahulu. Kita membangun ulang \omterm{def:b3:topology:continuity}{kekontinuan}, lalu menggarap
kekompakan lewat selimut terbuka (sambil membuktikan kesetaraannya dengan
definisi barisan pada Tahun ke-2 di ruang metrik), dan keterhubungan ---
lalu menutup dengan teorema bahwa bijeksi \omterm{def:b3:topology:continuity}{kontinu} tak dapat
menyamakan $\R$ dengan $\R^2$: jadi \omterm{def:b3:topology:topology}{topologi} dapat membedakan dimensi.

\section{Topologi, himpunan terbuka, kekontinuan}

\begin{definition}\label{def:b3:topology:topology}
Sebuah \emph{topologi}\index{topologi} pada himpunan $X$ adalah keluarga
$\mathcal T$ berisi himpunan bagian $X$ --- yang disebut \emph{himpunan
terbuka}\index{himpunan terbuka} --- sedemikian sehingga: $\varnothing, X \in
\mathcal T$; gabungan sembarang himpunan terbuka bersifat terbuka; dan irisan
\emph{berhingga} himpunan terbuka bersifat terbuka. Pasangan $(X, \mathcal T)$ disebut
\emph{ruang topologi}. Komplemen himpunan terbuka disebut
\emph{tertutup}. Adapun \emph{persekitaran}\index{persekitaran} sebuah $x$
adalah himpunan yang memuat suatu himpunan terbuka yang memuat $x$.
\end{definition}

\begin{example}\label{ex:b3:topology:examples}
(a) Ruang metrik, dengan ``terbuka'' seperti pada Tahun ke-2 (yaitu gabungan bola
terbuka): itulah \emph{\omterm{def:b3:topology:topology}{topologi} metrik}; dan metrik yang berbeda dapat memberikan
\omterm{def:b3:topology:topology}{topologi} yang sama (metrik ekuivalen). Ruang yang \omterm{def:b3:topology:topology}{topologinya}
lahir dari suatu metrik disebut \emph{dapat dimetrikan}.
(b) \omterm{def:b3:topology:topology}{Topologi} \emph{diskret} (semua himpunan bagiannya) dan
\omterm{def:b3:topology:topology}{topologi} \emph{indiskret} $\{\varnothing, X\}$.
(c) \omterm{def:b3:topology:topology}{Topologi} \emph{kofinit} pada himpunan tak berhingga: terbuka $=$
kosong atau berkomplemen berhingga. Ia tak dapat dimetrikan, seperti akan kita lihat
(\cref{exo:b3:topology:3}).
(d) Pada $\R$, \omterm{def:b3:topology:topology}{topologi} biasanya; pada $\bar\R = \R \cup
\{\pm\infty\}$, \omterm{def:b3:topology:topology}{topologi} urutan yang dibangkitkan sinar --- yang menjadikan
``$x_n \to +\infty$'' sekadar contoh kekonvergenan biasa.
\end{example}

\begin{definition}\label{def:b3:topology:interior}
Untuk $A \subseteq X$: \emph{interior} $\mathring A$ adalah
\omterm{def:b3:topology:topology}{himpunan terbuka} terbesar di dalam $A$ (gabungan semuanya);
\emph{penutup}\index{penutup} $\bar A$ adalah himpunan tertutup terkecil
yang memuat $A$; dan \emph{batas} $\partial A = \bar A
\setminus \mathring A$. Himpunan $A$ disebut \emph{padat}\index{himpunan padat} jika
$\bar A = X$. Berlaku $x \in \bar A$ jika dan hanya jika setiap \omterm{def:b3:topology:topology}{persekitaran}
$x$ memotong $A$ (sebab jika suatu \omterm{def:b3:topology:topology}{persekitaran} meleset dari $A$, komplemen inti
terbukanya menjadi himpunan tertutup yang lebih kecil di sekitar $A$; dan sebaliknya).
\end{definition}

\begin{definition}\label{def:b3:topology:basis}
Sebuah \emph{basis}\index{basis topologi} bagi $\mathcal T$ adalah
keluarga $\mathcal B \subseteq \mathcal T$ sedemikian sehingga setiap \omterm{def:b3:topology:topology}{himpunan
terbuka} merupakan gabungan anggota $\mathcal B$ (misalnya bola terbuka di
ruang metrik; selang terbuka di $\R$). Sebuah keluarga $\mathcal B$ berisi
himpunan bagian $X$ merupakan basis bagi \emph{suatu} \omterm{def:b3:topology:topology}{topologi} jika dan hanya jika ia menyelimuti
$X$ dan, untuk $B_1, B_2 \in \mathcal B$ serta $x \in B_1 \cap B_2$,
ada $B_3 \in \mathcal B$ dengan $x \in B_3 \subseteq B_1\cap B_2$
--- lalu ``gabungan anggotanya'' menjadi sebuah \omterm{def:b3:topology:topology}{topologi}, yaitu \omterm{def:b3:topology:topology}{topologi}
yang \emph{dibangkitkan} oleh $\mathcal B$.
\end{definition}

\begin{definition}\label{def:b3:topology:continuity}
Pemetaan $f \colon X \to Y$ disebut \emph{kontinu}\index{pemetaan kontinu} jika
$f^{-1}(V)$ terbuka untuk setiap $V \subseteq Y$ yang terbuka ---
setara dengan: prapeta himpunan tertutup bersifat tertutup; setara pula dengan:
untuk setiap $x$ dan \omterm{def:b3:topology:topology}{persekitaran} $V$ dari $f(x)$, $f^{-1}(V)$ merupakan
\omterm{def:b3:topology:topology}{persekitaran} $x$ (yakni kekontinuan \emph{di} setiap $x$). Cukup
diperiksa pada sebuah \omterm{def:b3:topology:basis}{basis} $Y$. Komposisi pemetaan kontinu
tetap kontinu. Sebuah \emph{homeomorfisma}\index{homeomorfisma}
adalah bijeksi kontinu yang inversnya kontinu; dan \omterm{def:b3:topology:topology}{topologi}
mempelajari sifat yang terpelihara oleh homeomorfisma.
\end{definition}

\begin{proposition}[Kekontinuan lawan penutup; perekatan]
\label{prop:b3:topology:gluing}
(a) Pemetaan $f$ \omterm{def:b3:topology:continuity}{kontinu} jika dan hanya jika $f(\bar A) \subseteq \overline{f(A)}$
untuk setiap $A \subseteq X$.
(b) Jika $X = F_1 \cup F_2$ dengan $F_1, F_2$ tertutup dan $f\colon X
\to Y$ terbatas secara \omterm{def:b3:topology:continuity}{kontinu} pada masing-masing $F_i$, maka $f$
\omterm{def:b3:topology:continuity}{kontinu}.
\end{proposition}

\begin{proof}
(a) Jika $f$ \omterm{def:b3:topology:continuity}{kontinu}: $f^{-1}(\overline{f(A)})$ tertutup dan
memuat $A$, sehingga memuat $\bar A$. Sebaliknya terapkan
kriterianya pada $A = f^{-1}(F)$ dengan $F$ tertutup: $f(\bar A) \subseteq
\overline{f(f^{-1}(F))} \subseteq \bar F = F$, sehingga $\bar A
\subseteq f^{-1}(F) = A$: jadi prapeta himpunan tertutup bersifat tertutup.
(b) Untuk $F \subseteq Y$ yang tertutup: $f^{-1}(F) =
(f\restriction_{F_1})^{-1}(F) \cup (f\restriction_{F_2})^{-1}
(F)$, yaitu gabungan dua himpunan yang tertutup di $F_1$ dan di $F_2$, sehingga
tertutup di $X$ (sebab $F_i$ tertutup: tertutup di dalam tertutup tetaplah tertutup).
\end{proof}

\begin{definition}\label{def:b3:topology:hausdorff}
Ruang $X$ disebut \emph{Hausdorff}\index{ruang Hausdorff} (atau
\emph{terpisah}) jika setiap dua titik yang berbeda mempunyai \omterm{def:b3:topology:topology}{persekitaran} yang
saling lepas. Ruang metrik bersifat Hausdorff (ambil bola berjari-jari
$d(x,y)/2$). Barisan $(x_n)$ \emph{konvergen} ke $x$ jika setiap
\omterm{def:b3:topology:topology}{persekitaran} $x$ memuat semua $x_n$ kecuali berhingga banyaknya; dan di ruang
Hausdorff limitnya tunggal (sebab dua limit akan mempunyai
\omterm{def:b3:topology:topology}{persekitaran} saling lepas yang masing-masing memuat sebuah ekor). Di ruang
Hausdorff, titik --- dan karenanya himpunan berhingga --- bersifat tertutup.
\end{definition}

\begin{remark}\label{rem:b3:topology:sequences}
Di ruang metrik, barisan mendeteksi segalanya: $x \in \bar A$
jika dan hanya jika ada barisan di $A$ yang konvergen ke $x$ (ambil $x_n \in A \cap
B(x, 1/n)$), dan $f$ \omterm{def:b3:topology:continuity}{kontinu} jika dan hanya jika ia \omterm{def:b3:topology:continuity}{kontinu} barisan.
Pada ruang umum kedua kesetaraan itu gagal; pengganti barisan
yang tepat (filter, jaring) termasuk kuliah yang lebih
lanjut. Kita akan menyatakan hasil berbasis barisan di ruang metrik
dan hasil berbasis selimut secara umum --- lalu membuktikan keduanya
setara di tempat mereka memang setara.
\end{remark}

\section{Subruang, hasil kali, kuosien}

\begin{definition}\label{def:b3:topology:constructions}
Ada tiga cara membuat ruang baru dari yang lama:\index{topologi
hasil kali}\index{topologi kuosien}
\begin{enumerate}
  \item \emph{Subruang}: pada $A \subseteq X$, \omterm{def:b3:topology:topology}{himpunan terbukanya} adalah
        $U \cap A$ dengan $U$ terbuka di $X$ --- yaitu \omterm{def:b3:topology:topology}{topologi}
        terkasar yang membuat inklusinya \omterm{def:b3:topology:continuity}{kontinu}.
  \item \emph{Hasil kali}: pada $X \times Y$ (dan hasil kali berhingga),
        yaitu \omterm{def:b3:topology:topology}{topologi} berbasis \emph{kotak terbuka} $U \times
        V$; sedangkan pada hasil kali tak berhingga $\prod_i X_i$, basisnya
        terdiri atas kotak $\prod U_i$ dengan $U_i = X_i$ untuk
        \emph{semua kecuali berhingga banyak} $i$ --- yaitu \omterm{def:b3:topology:topology}{topologi}
        terkasar yang membuat setiap proyeksinya \omterm{def:b3:topology:continuity}{kontinu}.
  \item \emph{Faktor}: jika $\sim$ sebuah relasi ekuivalensi pada $X$ dan
        $\pi\colon X \to X/{\sim}$ proyeksinya, nyatakan $V
        \subseteq X/{\sim}$ terbuka jika dan hanya jika $\pi^{-1}(V)$ terbuka ---
        yaitu \omterm{def:b3:topology:topology}{topologi} terhalus yang membuat $\pi$ \omterm{def:b3:topology:continuity}{kontinu}.
\end{enumerate}
\end{definition}

\begin{proposition}[Sifat universal]
\label{prop:b3:topology:universal}
(a) Pemetaan $f \colon Z \to X \times Y$ \omterm{def:b3:topology:continuity}{kontinu} jika dan hanya jika kedua
komponennya, $\pi_X \circ f$ dan $\pi_Y \circ f$, \omterm{def:b3:topology:continuity}{kontinu} (demikian pula untuk
hasil kali sembarang).
(b) Pemetaan $g \colon X/{\sim} \to Z$ \omterm{def:b3:topology:continuity}{kontinu} jika dan hanya jika $g \circ \pi
\colon X \to Z$ \omterm{def:b3:topology:continuity}{kontinu}.
\end{proposition}

\begin{proof}
(a) Keperluannya: lewat komposisi. Kecukupannya: cukup diperiksa
prapeta kotak basisnya: $f^{-1}(U \times V) = (\pi_Xf)^{-1}(U)
\cap (\pi_Yf)^{-1}(V)$, yang terbuka (sebab hanya berhingga faktor yang $\neq X_i$
pada kasus tak berhingga). (b) Keperluannya: lewat komposisi. Kecukupannya:
untuk $W \subseteq Z$ terbuka, $\pi^{-1}(g^{-1}(W)) = (g\pi)^{-1}(W)$
terbuka, dan menurut definisi \omterm{def:b3:topology:constructions}{topologi kuosiennya} itu berarti
$g^{-1}(W)$ terbuka.
\end{proof}

\begin{example}\label{ex:b3:topology:circle}
Kuosien $\R/\Z$ (dengan mengenali $x$ dan $x + n$) \omterm{def:b3:topology:continuity}{homeomorfik}
dengan lingkaran $S^1 = \{z \in \C : \abs z = 1\}$: sebab pemetaan $x
\mapsto \eu^{2\iu\pi x}$ beralih menjadi bijeksi \omterm{def:b3:topology:continuity}{kontinu}
$\R/\Z \to S^1$ (\cref{prop:b3:topology:universal}(b)); sedangkan
inversnya \omterm{def:b3:topology:continuity}{kontinu} menurut argumen kekompakan pada
\cref{cor:b3:topology:compacthomeo} di bawah
(\cref{exo:b3:topology:5} merincikan semuanya, termasuk mengapa
$\R/\Z$ bersifat \omterm{def:b3:topology:hausdorff}{Hausdorff} dan \omterm{def:b3:topology:compact}{kompak}). Demikian pula $[0,1]$ dengan
kedua ujungnya direkatkan adalah $S^1$, persegi dengan sisi berhadapannya direkatkan
adalah torus, dan perekatan akhirnya menjadi teorema, bukan gambar.
\end{example}

\section{Kekompakan}

\begin{definition}\label{def:b3:topology:compact}
Sebuah \emph{selimut terbuka} bagi $X$ adalah keluarga $(U_i)_{i\in I}$ berisi \omterm{def:b3:topology:topology}{himpunan
terbuka} dengan $\bigcup U_i = X$. Ruang $X$ disebut
\emph{kompak}\index{ruang kompak} jika ia \omterm{def:b3:topology:hausdorff}{Hausdorff} dan setiap
selimut terbukanya mempunyai \emph{subselimut berhingga}. Setara dengan itu (dengan mengambil
komplemennya): setiap keluarga himpunan tertutup yang bersifat \emph{irisan
berhingga} (yakni semua subkeluarga berhingganya beririsan tak kosong)
mempunyai irisan total yang tak kosong.
\end{definition}

\begin{theorem}[Sifat pertama]\label{thm:b3:topology:compactprops}
Misalkan $X$ \omterm{def:b3:topology:compact}{kompak}.
\begin{enumerate}
  \item Himpunan bagian tertutup dari $X$ bersifat \omterm{def:b3:topology:compact}{kompak}; dan himpunan bagian \omterm{def:b3:topology:compact}{kompak} dari
        \omterm{def:b3:topology:hausdorff}{ruang Hausdorff} bersifat tertutup.
  \item Peta \omterm{def:b3:topology:continuity}{kontinu} \omterm{def:b3:topology:compact}{ruang kompak} di dalam \omterm{def:b3:topology:hausdorff}{ruang Hausdorff}
        bersifat \omterm{def:b3:topology:compact}{kompak}. Khususnya $f\colon X
        \to \R$ yang \omterm{def:b3:topology:continuity}{kontinu} bersifat terbatas dan mencapai batasnya.
  \item Barisan turun himpunan bagian tertutup tak kosong dari $X$
        mempunyai irisan yang tak kosong.
\end{enumerate}
\end{theorem}

\begin{proof}
(1) Misalkan $F \subseteq X$ tertutup dan $(U_i)$ sebuah selimut terbuka bagi $F$
(dengan \omterm{def:b3:topology:topology}{himpunan terbuka} $X$): menambahkan $X \setminus F$ memberikan selimut terbuka bagi
$X$; lalu subselimut berhingganya, setelah $X \setminus F$ dibuang, menyelimuti $F$.
Sifat \omterm{def:b3:topology:hausdorff}{Hausdorff} pada subruangnya jelas. Sebaliknya misalkan $K \subseteq Y$
\omterm{def:b3:topology:compact}{kompak}, $Y$ \omterm{def:b3:topology:hausdorff}{Hausdorff}, dan $y \notin K$: untuk setiap $x \in K$
pilihlah \omterm{def:b3:topology:topology}{himpunan terbuka} saling lepas $U_x \ni x$, $V_x \ni y$; lalu berhingga banyak $U_x$
menyelimuti $K$, dan irisan $V_x$ yang bersesuaian menjadi
\omterm{def:b3:topology:topology}{persekitaran} $y$ yang saling lepas dengan selimut $K$, sehingga saling lepas dengan $K$:
jadi komplemen $K$ terbuka.

(2) Jika $(V_j)$ menyelimuti $f(X)$, maka $(f^{-1}(V_j))$ menyelimuti $X$; dan
subselimut berhingga di bawah berasal dari subselimut berhingga di
atas. Petanya bersifat \omterm{def:b3:topology:hausdorff}{Hausdorff} sebagai subruang. Untuk $f$ yang real:
$f(X)$ \omterm{def:b3:topology:compact}{kompak} di $\R$, sehingga tertutup dan terbatas (selimuti dengan
$(-n, n)$ untuk keterbatasannya; dan tertutup menurut (1)), sedangkan himpunan tertutup terbatas
memuat supremumnya.

(3) Jika $\bigcap F_n = \varnothing$, maka \omterm{def:b3:topology:topology}{himpunan terbuka} $X \setminus F_n$
menyelimuti $X$; berhingga banyak di antaranya sudah cukup, sehingga ada $F_{n_0} =
\varnothing$ (karena barisannya turun) --- kontradiksi.
\end{proof}

\begin{corollary}\label{cor:b3:topology:compacthomeo}
Sebuah \emph{bijeksi} \omterm{def:b3:topology:continuity}{kontinu} dari \omterm{def:b3:topology:compact}{ruang kompak} ke \omterm{def:b3:topology:hausdorff}{ruang
Hausdorff} merupakan \omterm{def:b3:topology:continuity}{homeomorfisma}.
\end{corollary}

\begin{proof}
Inversnya \omterm{def:b3:topology:continuity}{kontinu} jika dan hanya jika peta langsung himpunan tertutup bersifat
tertutup; sedangkan $F$ yang tertutup bersifat \omterm{def:b3:topology:compact}{kompak}
(\cref{thm:b3:topology:compactprops}(1)), petanya \omterm{def:b3:topology:compact}{kompak} menurut
(2), sehingga tertutup menurut (1) di dalam sasaran \omterm{def:b3:topology:hausdorff}{Hausdorffnya}.
\end{proof}

\begin{theorem}[Hasil kali berhingga]\label{thm:b3:topology:tychonoff}
Hasil kali berhingga \omterm{def:b3:topology:compact}{ruang kompak} bersifat \omterm{def:b3:topology:compact}{kompak}.\index{teorema Tychonoff
(kasus berhingga)}
\end{theorem}

\begin{proof}
Cukup ditangani $X \times Y$. Sifat \omterm{def:b3:topology:hausdorff}{Hausdorffnya} diwarisi
(pisahkan pada satu koordinat). Misalkan $(W_i)$ sebuah selimut terbuka bagi $X
\times Y$; kita boleh menganggap $W_i$ berupa kotak \omterm{def:b3:topology:basis}{basis} $U_i \times
V_i$ (haluskan saja: setiap titik duduk di sebuah kotak di dalam suatu $W_i$;
lalu subselimut berhingga atas kotaknya melahirkan subselimut berhingga atas $W_i$). Tetapkan $x \in X$:
maka \emph{keping} $\{x\}\times Y \cong Y$ bersifat \omterm{def:b3:topology:compact}{kompak}, sehingga berhingga
banyak kotak $U_1\times V_1, \dots, U_k\times V_k$ menyelimutinya, dengan
$x \in U_j$ untuk setiap $j$; lalu $U^x = \bigcap_{j \leq k} U_j$ menjadi
\omterm{def:b3:topology:topology}{persekitaran} terbuka $x$ dengan $U^x \times Y$ terselimuti oleh
berhingga kotak (inilah \emph{lema tabung}\index{lema tabung}:
sebab untuk $(x', y) \in U^x \times Y$ berlaku $y \in V_j$
untuk suatu $j$ karena kotak itu menyelimuti $\{x\}\times Y$ pada ketinggian
$y$, dan $x' \in U^x \subseteq U_j$, sehingga $(x', y) \in U_j
\times V_j$). Kini berhingga banyak $U^x$
menyelimuti $X$ yang \omterm{def:b3:topology:compact}{kompak}; dan kumpulan berhingga kotak yang bersesuaian
menyelimuti $X \times Y$.
\end{proof}

\begin{theorem}[Kekompakan di ruang metrik]
\label{thm:b3:topology:metriccompact}
Untuk sebuah ruang metrik $(X, d)$, berikut ini saling
setara:\index{kekompakan barisan}
\begin{enumerate}
  \item $X$ \omterm{def:b3:topology:compact}{kompak} (Borel--Lebesgue);
  \item setiap barisan di $X$ mempunyai subbarisan yang konvergen
        (\omterm{thm:b3:topology:metriccompact}{kekompakan barisan} --- definisi Tahun ke-2);
  \item $X$ lengkap dan \emph{terbatas total}: yakni untuk setiap
        $\varepsilon > 0$, berhingga banyak bola berjari-jari
        $\varepsilon$ menyelimuti $X$.
\end{enumerate}
\end{theorem}

\begin{proof}
(1)$\Rightarrow$(2): Misalkan $(x_n)$ tak punya subbarisan yang
konvergen. Maka setiap $x \in X$ mempunyai bola terbuka $B_x$ yang
memuat $x_n$ hanya untuk berhingga banyak $n$ (sebab kalau tidak, ada
subbarisan yang konvergen ke $x$: ambil jari-jari $1/k$). Berhingga banyak
$B_x$ menyelimuti $X$, sehingga hanya ada berhingga indeks $n$:
mustahil.

(2)$\Rightarrow$(3): Kelengkapannya: barisan Cauchy yang mempunyai
subbarisan konvergen akan konvergen (Tahun ke-2). Keterbatasan totalnya: jika
ada $\varepsilon$ yang tak punya selimut berhingga, pilihlah secara induktif
$x_{n+1}$ di luar $\bigcup_{k\leq n} B(x_k, \varepsilon)$: maka
barisannya memenuhi $d(x_m, x_n) \geq \varepsilon$ untuk $m \neq n$, sehingga tak punya
subbarisan Cauchy dan tak punya subbarisan konvergen.

(3)$\Rightarrow$(1): Pertama, (3) mengakibatkan (2): diberikan $(x_n)$,
selimuti $X$ dengan berhingga bola berjari-jari $1$: salah satunya, sebut $B_1$,
memuat sebuah subbarisan; lalu selimuti $X$ dengan bola berjari-jari $1/2$: salah satunya
memuat subbarisan lanjutannya; iterasikan lalu diagonalkan: subbarisan
diagonalnya bersifat Cauchy (karena dua sukunya setelah tahap $k$ terletak
di satu bola berjari-jari $2^{-k}$, sampai pada faktor $2\cdot$ yang biasa),
sehingga konvergen. Sekarang misalkan $(U_i)$ sebuah selimut terbuka dan andaikan tak ada
subselimut berhingga. \emph{Argumen bilangan Lebesgue}: untuk setiap $n$,
ada bola $B(y_n, 2^{-n})$ yang tak terselimuti berhingga banyak
$U_i$ --- memang, selimuti $X$ dengan berhingga bola berjari-jari
$2^{-n}$; seandainya masing-masing terselimuti secara berhingga, maka $X$ pun demikian. Menurut
(2), ada subbarisan $y_{n_k} \to y$; pilihlah $i$ dengan $y \in U_i$
dan $r > 0$ dengan $B(y, r) \subseteq U_i$. Untuk $k$ besar,
$B(y_{n_k}, 2^{-n_k}) \subseteq B(y, r) \subseteq U_i$: jadi terselimuti
oleh \emph{satu} $U_i$ --- kontradiksi.
\end{proof}

\begin{corollary}[Heine--Borel; Heine]
\label{cor:b3:topology:heineborel}
(a) Himpunan bagian $\R^n$ bersifat \omterm{def:b3:topology:compact}{kompak} jika dan hanya jika ia tertutup dan terbatas.
(b) \omterm{def:b3:topology:continuity}{Pemetaan kontinu} dari ruang metrik \omterm{def:b3:topology:compact}{kompak} ke ruang
metrik bersifat \omterm{def:b3:topology:continuity}{kontinu} seragam.\index{teorema Heine--Borel}\index{teorema Heine}
\end{corollary}

\begin{proof}
(a) Tertutup dan terbatas $\Rightarrow$ termuat di sebuah kubus
$[-M,M]^n$, yang bersifat \omterm{def:b3:topology:compact}{kompak}: sebab $[-M, M]$ demikian (secara barisan, menurut
Bolzano--Weierstrass --- atau langsung lewat dikotomi pada selimutnya),
dan \cref{thm:b3:topology:tychonoff} menangani hasil kalinya; lalu
terapkan \cref{thm:b3:topology:compactprops}(1). Sebaliknya himpunan bagian
\omterm{def:b3:topology:compact}{kompak} bersifat tertutup
(\cref{thm:b3:topology:compactprops}(1)) dan terbatas (selimuti dengan
bola sepusat).

(b) Misalkan $f \colon X \to Y$ dan $\varepsilon > 0$. Bola
$B\bigl(x, \delta_x\bigr)$ dengan $f\bigl(B(x,
2\delta_x)\bigr)\subseteq B\bigl(f(x), \varepsilon/2\bigr)$
menyelimuti $X$; ekstraklah subselimut berhingga $B(x_j, \delta_{x_j})$ lalu
tetapkan $\delta = \min_j \delta_{x_j}$. Jika $d(x, x') < \delta$: maka $x
\in B(x_j, \delta_{x_j})$ untuk suatu $j$, dan lalu keduanya memenuhi $x, x'
\in B(x_j, 2\delta_{x_j})$, sehingga $d(f(x), f(x')) < \varepsilon$.
\end{proof}

\begin{definition}\label{def:b3:topology:locallycompact}
Ruang $X$ disebut \emph{kompak lokal}\index{ruang kompak lokal} jika ia
\omterm{def:b3:topology:hausdorff}{Hausdorff} dan setiap titiknya mempunyai \omterm{def:b3:topology:topology}{persekitaran} \omterm{def:b3:topology:compact}{kompak} (misalnya $\R^n$;
himpunan bagian terbuka $\R^n$; ruang diskret --- tetapi bukan $\Q$, lihat
\cref{exo:b3:topology:9}). Setiap \omterm{def:b3:topology:compact}{ruang kompak} lokal terbenam di dalam
\omterm{def:b3:topology:compact}{ruang kompak}: yaitu \emph{kompaktifikasi satu
titik}\index{kompaktifikasi satu titik} $\hat X = X
\cup \{\infty\}$, yang \omterm{def:b3:topology:topology}{himpunan terbukanya} adalah \omterm{def:b3:topology:topology}{himpunan terbuka} $X$ beserta
komplemen (di $\hat X$) himpunan bagian \omterm{def:b3:topology:compact}{kompak} $X$. Aksiomanya
mudah diperiksa langsung; lalu $\hat X$ bersifat \omterm{def:b3:topology:compact}{kompak} (sebab selimutnya mempunyai anggota
yang memuat $\infty$, yang komplemennya \omterm{def:b3:topology:compact}{kompak} dan terselimuti oleh
berhingga anggota lain) sekaligus \omterm{def:b3:topology:hausdorff}{Hausdorff} (pisahkan $x$ dari $\infty$
lewat \omterm{def:b3:topology:topology}{persekitaran} \omterm{def:b3:topology:compact}{kompak} $x$ dan komplemennya). Contohnya:
$\hat\R \cong S^1$, dan $\widehat{\R^n} \cong S^n$ lewat
proyeksi stereografis (\cref{exo:b3:topology:11}).
\end{definition}

\section{Keterhubungan}

\begin{definition}\label{def:b3:topology:connected}
Ruang $X$ disebut \emph{terhubung}\index{ruang terhubung} jika ia bukan
gabungan dua \omterm{def:b3:topology:topology}{himpunan terbuka} tak kosong yang saling lepas --- setara dengan:
satu-satunya himpunan bagiannya yang sekaligus terbuka dan tertutup adalah $\varnothing$ dan $X$;
setara pula dengan: setiap \omterm{def:b3:topology:continuity}{pemetaan kontinu} $X \to \{0, 1\}$ (diskret)
bersifat konstan. Himpunan bagian disebut terhubung jika ia terhubung sebagai subruang.
\end{definition}

\begin{theorem}\label{thm:b3:topology:connectedbasics}
\begin{enumerate}
  \item Himpunan bagian $\R$ yang \omterm{def:b3:topology:connected}{terhubung} tepat berupa selang.
  \item Peta \omterm{def:b3:topology:continuity}{kontinu} himpunan \omterm{def:b3:topology:connected}{terhubung} tetap \omterm{def:b3:topology:connected}{terhubung} (dan dari sini
        teorema nilai antara: pemetaan real \omterm{def:b3:topology:continuity}{kontinu} pada
        \omterm{def:b3:topology:connected}{ruang terhubung} berpeta sebuah selang).
  \item Jika $(A_i)$ \omterm{def:b3:topology:connected}{terhubung} dan mempunyai titik persekutuan, maka $\bigcup
        A_i$ \omterm{def:b3:topology:connected}{terhubung}. Jika $A$ \omterm{def:b3:topology:connected}{terhubung} dan $A \subseteq
        B \subseteq \bar A$, maka $B$ \omterm{def:b3:topology:connected}{terhubung}.
  \item Hasil kali berhingga \omterm{def:b3:topology:connected}{ruang terhubung} bersifat \omterm{def:b3:topology:connected}{terhubung}.
\end{enumerate}
\end{theorem}

\begin{proof}
Sepanjang bukti ini kita memakai kriteria $\{0,1\}$: bahwa $f \colon X \to
\{0,1\}$ yang \omterm{def:b3:topology:continuity}{kontinu} haruslah konstan.

(1) Himpunan $A$ yang bukan selang meleset dari suatu $c$ di antara $a, b \in A$:
lalu $A = (A \cap (-\infty, c)) \sqcup (A \cap (c, +\infty))$
memutuskannya. Sebaliknya misalkan $I$ sebuah selang dan $f \colon I
\to \{0,1\}$ \omterm{def:b3:topology:continuity}{kontinu} dengan $f(a) = 0$, $f(b) = 1$, $a < b$.
Misalkan $c = \sup\{x \in [a,b] : f(x) = 0\}$; \omterm{def:b3:topology:continuity}{kekontinuan} di $c$
memaksa $f(c) = 0$ (sebab ia limit nilai $0$: setiap \omterm{def:b3:topology:topology}{persekitaran}
$c$ memotong $\{f = 0\}$, dan $\{f=0\}$ tertutup) lalu $c < b$
dengan $f \equiv 1$ pada $(c, b]$, sehingga $f(c) = 1$ menurut
argumen \omterm{def:b3:topology:interior}{penutup} yang sama pada $\{f = 1\} \ni c$: kontradiksi.

(2) \omterm{def:b3:topology:continuity}{Pemetaan kontinu} $g \colon f(X) \to \{0,1\}$ melahirkan $g \circ f$
yang konstan, sehingga $g$ konstan pada petanya.

(3) \omterm{def:b3:topology:continuity}{Pemetaan kontinu} $f\colon \bigcup A_i \to \{0,1\}$ bersifat konstan
pada setiap $A_i$, dengan nilai yang sama di titik persekutuannya. Untuk
\omterm{def:b3:topology:interior}{penutupnya}: $f \colon B \to \{0,1\}$ konstan $= c$ pada $A$; dan sembarang
$x \in B \subseteq \bar A$ berada di \omterm{def:b3:topology:interior}{penutup} $A$, sedangkan $f^{-1}
(f(x))$ merupakan \omterm{def:b3:topology:topology}{persekitaran} $x$ (sebab prapeta \omterm{def:b3:topology:topology}{himpunan terbuka}), yang tentu
memotong $A$: jadi $f(x) = c$.

(4) Untuk $X \times Y$ dan $f\colon X\times Y \to \{0,1\}$: sembarang
dua titik $(x,y), (x',y')$ dihubungkan oleh ``siku''
$\{x\}\times Y \cup X \times \{y'\}$, yaitu gabungan dua himpunan \omterm{def:b3:topology:connected}{terhubung}
(yang \omterm{def:b3:topology:continuity}{homeomorfik} dengan $Y$ dan $X$) yang bertemu di $(x, y')$: jadi menurut (3) dan
(2), $f$ bernilai sama di kedua titik itu.
\end{proof}

\begin{definition}\label{def:b3:topology:pathconnected}
Ruang $X$ disebut \emph{terhubung lintasan}\index{ruang terhubung lintasan} jika sembarang
dua titiknya dihubungkan oleh sebuah \emph{lintasan} (yakni $\gamma
\colon [0,1] \to X$ yang \omterm{def:b3:topology:continuity}{kontinu}). \omterm{def:b3:topology:connected}{Terhubung} lintasan mengakibatkan \omterm{def:b3:topology:connected}{terhubung}: sebab dua
nilai \omterm{def:b3:topology:continuity}{pemetaan kontinu} $f \colon X \to \{0,1\}$ di $x, y$ merupakan
nilai $f\circ\gamma$ yang konstan
(\cref{thm:b3:topology:connectedbasics}(1)--(2)). Himpunan bagian konveks
ruang bernorma bersifat \omterm{def:b3:topology:connected}{terhubung} lintasan (yaitu lewat ruas garis); demikian pula $\R^n
\setminus \{0\}$ untuk $n \geq 2$ (kelilingi titik asalnya), dan
$S^n$ untuk $n \geq 1$ (proyeksikan lintasan dari $\R^{n+1}\setminus
\{0\}$).
\end{definition}

\begin{example}[Kurva sinus topolog]
\label{ex:b3:topology:sinecurve}
Misalkan $\Gamma = \{(x, \sin\frac1x) : 0 < x \leq 1\}$ dan $S =
\bar\Gamma = \Gamma \cup (\{0\}\times[-1,1])$ (sebab setiap titik $(0,
y)$ dengan $\abs y \leq 1$ merupakan limit titik $\Gamma$: selesaikan
$\sin\frac1x = y$ di dekat $0$). Maka $S$ bersifat
\emph{\omterm{def:b3:topology:connected}{terhubung}} --- sebagai \omterm{def:b3:topology:interior}{penutup} $\Gamma$ yang \omterm{def:b3:topology:connected}{terhubung}, yaitu
peta \omterm{def:b3:topology:continuity}{kontinu} $(0, 1]$
(\cref{thm:b3:topology:connectedbasics}(3)) --- tetapi \emph{tidak
\omterm{def:b3:topology:pathconnected}{terhubung lintasan}}: sebab \omterm{def:b3:topology:pathconnected}{lintasan} dari $(1, \sin 1)$ ke $(0,0)$ harus
melintasi absis yang $\to 0$ sementara ordinatnya berayun
antara $\pm1$; \cref{exo:b3:topology:9} membuatnya ketat.
Keterhubungan dan keterhubungan \omterm{def:b3:topology:pathconnected}{lintasan} memang benar-benar
berbeda.\index{kurva sinus topolog}
\end{example}

\begin{omfigure}
\begin{tikzpicture}
\begin{axis}[omaxis, width=10.5cm, height=4.6cm,
    xmin=-0.02, xmax=0.55, ymin=-1.25, ymax=1.25,
    xtick={0, 0.25, 0.5}, ytick={-1, 1}]
  \addplot[omThm, thick, domain=0.006:0.55, samples=1200]
    {sin(deg(1/x))};
  \addplot[omDef, very thick] coordinates {(0,-1) (0,1)};
\end{axis}
\end{tikzpicture}

{\small \omterm{ex:b3:topology:sinecurve}{Kurva sinus topolog}: grafik
$\sin\frac1x$ (merah) menumpuk pada seluruh ruas garis
$\{0\}\times[-1,1]$ (biru). Gabungannya \omterm{def:b3:topology:connected}{terhubung} tetapi tidak
\omterm{def:b3:topology:pathconnected}{terhubung lintasan}: tak ada \omterm{def:b3:topology:pathconnected}{lintasan} \omterm{def:b3:topology:continuity}{kontinu} yang dapat menyeberangi ayunan
yang tak berhingga banyaknya itu dalam waktu parameter yang berhingga.}
\end{omfigure}

\begin{definition}\label{def:b3:topology:components}
Yang disebut \emph{komponen terhubung}\index{komponen terhubung} dari $x
\in X$ adalah gabungan semua himpunan bagian \omterm{def:b3:topology:connected}{terhubung} yang memuat $x$ ---
yakni yang terbesar (\cref{thm:b3:topology:connectedbasics}(3)).
Komponen mempartisi $X$ dan bersifat tertutup (sebab \omterm{def:b3:topology:interior}{penutup} himpunan
\omterm{def:b3:topology:connected}{terhubung} tetap \omterm{def:b3:topology:connected}{terhubung}). Ruang $X$ disebut \emph{tak \omterm{def:b3:topology:connected}{terhubung} total} jika semua
komponennya beranggota tunggal (misalnya $\Q$; atau himpunan Cantor pada soal akhir
pekan).
\end{definition}

\begin{theorem}\label{thm:b3:topology:rnotr2}
$\R$ tidak \omterm{def:b3:topology:continuity}{homeomorfik} dengan $\R^n$ mana pun yang $n \geq 2$.
\end{theorem}

\begin{proof}
Andaikan $h \colon \R \to \R^n$ sebuah \omterm{def:b3:topology:continuity}{homeomorfisma}. Dengan membuang sebuah
titik: $\R \setminus \{0\}$ \omterm{def:b3:topology:continuity}{homeomorfik} dengan $\R^n \setminus
\{h(0)\}$. Padahal $\R\setminus\{0\}$ tak \omterm{def:b3:topology:connected}{terhubung}, sedangkan
$\R^n\setminus\{\text{titik}\}$ \omterm{def:b3:topology:pathconnected}{terhubung lintasan} untuk $n \geq
2$ (\cref{def:b3:topology:pathconnected}; geser titiknya ke
$0$): dan keterhubungan merupakan invarian \omterm{def:b3:topology:continuity}{homeomorfisma}
(\cref{thm:b3:topology:connectedbasics}(2)) --- kontradiksi.
(Bahwa $\R^2 \not\cong \R^3$ menuntut invarian yang lebih halus ---
yaitu \omterm{def:b3:topology:topology}{topologi} aljabar; dan soal akhir pekan menunjukkan bahayanya:
\emph{surjeksi \omterm{def:b3:topology:continuity}{kontinu}} $\R \to \R^2$ memang ada.)
\end{proof}

\begin{method}\label{met:b3:topology:connectargument}
Untuk membuktikan sebuah himpunan \omterm{def:b3:topology:connected}{terhubung}: tampilkan ia sebagai peta \omterm{def:b3:topology:continuity}{kontinu},
sebagai gabungan himpunan \omterm{def:b3:topology:connected}{terhubung} yang saling bertindih, sebagai \omterm{def:b3:topology:interior}{penutup}, atau sebagai hasil kali
(\cref{thm:b3:topology:connectedbasics}); dan untuk himpunan bagian ruang
bernorma, buktikan keterhubungan \omterm{def:b3:topology:pathconnected}{lintasannya} dengan \omterm{def:b3:topology:pathconnected}{lintasan} eksplisit (ruas garis,
busur, siku). Untuk membuktikan bahwa dua ruang \emph{tidak} \omterm{def:b3:topology:continuity}{homeomorfik}:
carilah invarian \omterm{def:b3:topology:topology}{topologi} yang berbeda --- kekompakan,
keterhubungan, banyaknya komponen, atau komponennya setelah
sebuah himpunan berhingga yang dipilih baik-baik dibuang (yaitu muslihat
\cref{thm:b3:topology:rnotr2}: ia juga membuktikan $[0,1) \not\cong
(0,1)$ dan $S^1 \not\cong [0,1]$).
\end{method}

\section{Latihan}

\begin{exercise}[$\star$]\label{exo:b3:topology:1}
(a) Daftarkan semua \omterm{def:b3:topology:topology}{topologi} pada $\{a, b\}$, klasifikasikan sampai
\omterm{def:b3:topology:continuity}{homeomorfisma}, lalu tentukan mana yang \omterm{def:b3:topology:connected}{terhubung} dan mana
yang \omterm{def:b3:topology:hausdorff}{Hausdorff}.
(b) Di $\R$ (dengan \omterm{def:b3:topology:topology}{topologi} biasa), hitunglah \omterm{def:b3:topology:interior}{interior}, \omterm{def:b3:topology:interior}{penutup} dan
batas dari $\Q$, dari $[0,1) \cup \{2\}$, dan dari $\{1/n : n \geq
1\}$.
\end{exercise}

\begin{exercise}[$\star$]\label{exo:b3:topology:2}
(a) Tunjukkan bahwa $f\colon X \to Y$ \omterm{def:b3:topology:continuity}{kontinu} jika dan hanya jika $f^{-1}(B)$
terbuka untuk setiap $B$ pada sebuah \omterm{def:b3:topology:basis}{basis} tetap $Y$.
(b) Tunjukkan bahwa $f = \mathbf 1_{[0,\infty)} \colon \R \to \R$
tidak \omterm{def:b3:topology:continuity}{kontinu} tetapi \emph{\omterm{def:b3:topology:continuity}{kontinu} kanan}. Periksalah bahwa
selang setengah terbuka $[a, b)$ membentuk \omterm{def:b3:topology:basis}{basis} sebuah \omterm{def:b3:topology:topology}{topologi} pada
$\R$ sebagai sumbernya (yaitu \emph{garis Sorgenfrey}), bahwa \omterm{def:b3:topology:topology}{topologi} itu
lebih halus secara sejati daripada \omterm{def:b3:topology:topology}{topologi} biasanya, dan bahwa fungsi $\R \to
\R$ (dengan sasaran biasa) \omterm{def:b3:topology:continuity}{kontinu} dari garis Sorgenfrey jika dan hanya jika ia
\omterm{def:b3:topology:continuity}{kontinu} kanan di setiap titik.
\end{exercise}

\begin{exercise}[$\star$]\label{exo:b3:topology:3}
Pada himpunan tak berhingga dengan \omterm{def:b3:topology:topology}{topologi} kofinit, tunjukkan: setiap dua
\omterm{def:b3:topology:topology}{himpunan terbuka} tak kosong saling memotong (sehingga ruangnya tidak \omterm{def:b3:topology:hausdorff}{Hausdorff}, dan karenanya
tidak dapat dimetrikan); dan setiap barisan injektif konvergen ke
\emph{setiap} titik. Di manakah bukti ketunggalan limit
memakai sifat \omterm{def:b3:topology:hausdorff}{Hausdorff}?
\end{exercise}

\begin{exercise}[$\star\star$]\label{exo:b3:topology:4}
(a) Tunjukkan bahwa proyeksi $\pi_X, \pi_Y$ sebuah hasil kali bersifat
\omterm{def:b3:topology:continuity}{kontinu} dan \emph{terbuka} (yakni peta \omterm{def:b3:topology:topology}{himpunan terbuka} tetap terbuka), tetapi umumnya tidak
tertutup (lihat $\{xy = 1\}$ di $\R^2$).
(b) Tunjukkan bahwa barisan di dalam hasil kali terbilang $\prod_n X_n$ atas
ruang metrik konvergen jika dan hanya jika setiap koordinatnya konvergen, dan bahwa
$d(x, y) = \sum_n 2^{-n}\min(1, d_n(x_n, y_n))$ memetrikan
\omterm{def:b3:topology:constructions}{topologi hasil kalinya}.
(c) Pada $\R^\N$, tunjukkan bahwa ``\omterm{def:b3:topology:topology}{topologi} kotak'' (yakni semua hasil kali
\omterm{def:b3:topology:topology}{himpunan terbuka} bersifat terbuka) lebih halus secara sejati: barisan
$x^{(k)} = (1/k, 1/k, \dots)$ konvergen ke $0$ pada \omterm{def:b3:topology:constructions}{topologi
hasil kalinya} tetapi tidak pada \omterm{def:b3:topology:topology}{topologi} kotaknya.
\end{exercise}

\begin{exercise}[$\star\star$]\label{exo:b3:topology:5}
Lingkaran, dengan tiga cara. Tunjukkan bahwa berikut ini saling
\omterm{def:b3:topology:continuity}{homeomorfik} berpasangan, disertai pemetaan eksplisitnya: (i) $S^1 \subseteq \R^2$;
(ii) $\R/\Z$ (dengan \omterm{def:b3:topology:constructions}{topologi kuosien}); (iii) $[0,1]/(0 \sim 1)$.
\emph{(Untuk (ii): tunjukkan bahwa $\R/\Z$ bersifat \omterm{def:b3:topology:hausdorff}{Hausdorff} --- angkatlah dua kelasnya
menjadi wakil yang berjarak $\leq \frac12$ --- dan bahwa
$\pi([0,1])$ adalah seluruhnya, sehingga kuosiennya \omterm{def:b3:topology:compact}{kompak}; lalu pakai
\cref{cor:b3:topology:compacthomeo}.)}
\end{exercise}

\begin{exercise}[$\star\star$]\label{exo:b3:topology:6}
Misalkan $X$ sebuah ruang metrik.
(a) Tunjukkan bahwa gabungan berhingga himpunan bagian \omterm{def:b3:topology:compact}{kompak} bersifat \omterm{def:b3:topology:compact}{kompak}, dan
bahwa irisan sembarangnya pun demikian.
(b) Jika $K$ \omterm{def:b3:topology:compact}{kompak}, $F$ tertutup, dan $K \cap F = \varnothing$,
tunjukkan $d(K, F) = \inf\{d(x,y) : x \in K, y \in F\} > 0$; lalu berilah
contoh penyangkal dengan dua himpunan tertutup yang saling lepas.
(c) Tunjukkan bahwa $X$ \omterm{def:b3:topology:compact}{kompak} jika dan hanya jika \emph{setiap} $f
\colon X \to \R$ yang \omterm{def:b3:topology:continuity}{kontinu} bersifat terbatas. \emph{(Jika ada barisan yang tak punya
subbarisan konvergen, bangunlah fungsi \omterm{def:b3:topology:continuity}{kontinu} tak terbatas yang
berpendukung dekat sukunya; atau pakailah fungsi bertipe $x \mapsto
d(x, \cdot)$ --- satu jalur yang bersih: jika $(x_n)$ tak punya
titik kumpul, maka himpunan $\{x_n\}$ tertutup dan diskret, dan
$f(x_n) = n$ dapat diperluas secara \omterm{def:b3:topology:continuity}{kontinu} tanpa Tietze lewat
$f(x) = \sum_n n\,\max\bigl(0, 1 - \frac{d(x, x_n)}{r_n}\bigr)$
untuk $r_n$ yang cukup kecil.)}
\end{exercise}

\begin{exercise}[$\star\star$]\label{exo:b3:topology:7}
(Bilangan Lebesgue) Misalkan $(U_i)$ sebuah selimut terbuka bagi ruang
metrik \omterm{def:b3:topology:compact}{kompak} $X$. Tunjukkan bahwa ada $\delta > 0$ sedemikian sehingga setiap
himpunan bagian berdiameter $< \delta$ terletak di dalam satu $U_i$ saja.
\emph{(Kalau tidak, pilihlah $A_n$ berdiameter $< 1/n$ yang tak termuat di satu $U_i$ pun, lalu
ambil titik kumpul dari $x_n \in A_n$ yang dipilih.)} Lalu simpulkan kembali teorema
Heine (\cref{cor:b3:topology:heineborel}(b)).
\end{exercise}

\begin{exercise}[$\star\star$]\label{exo:b3:topology:8}
(a) Tunjukkan bahwa $GL_n(\R)$ merupakan himpunan bagian $M_n(\R)$ yang terbuka dan \omterm{def:b3:topology:interior}{padat},
dan bahwa ia tak \omterm{def:b3:topology:connected}{terhubung}: sebab tanda determinannya
memisahkannya menjadi (sekurang-kurangnya) dua kepingan. Tunjukkan sebaliknya
bahwa $GL_n(\C)$ \omterm{def:b3:topology:pathconnected}{terhubung lintasan}.
\emph{(Untuk $A, B \in GL_n(\C)$: $\lambda \mapsto \det\bigl((1 -
\lambda)A + \lambda B\bigr)$ merupakan polinomial yang tak nol secara
identik, sehingga berakar berhingga banyak di $\C$: pilihlah \omterm{def:b3:topology:pathconnected}{lintasan}
$\lambda$ di $\C$ dari $0$ ke $1$ yang menghindari akar itu.)}
(b) Buktikan kepadatannya: $A + \varepsilon I$ terbalikkan untuk $\varepsilon > 0$
yang kecil.
\end{exercise}

\begin{exercise}[$\star\star$]\label{exo:b3:topology:9}
(a) Tunjukkan bahwa \omterm{def:b3:topology:components}{komponen terhubung} $\Q$ berupa himpunan
beranggota tunggal, dan bahwa $\Q$ tidak \omterm{def:b3:topology:locallycompact}{kompak lokal} (sebab \omterm{def:b3:topology:topology}{persekitaran} \omterm{def:b3:topology:compact}{kompak}
$0$ di $\Q$ akan memuat $[-r, r]\cap\Q$,
yang \omterm{def:b3:topology:interior}{penutupnya} di $\Q$ tidak \omterm{def:b3:topology:compact}{kompak}: potonglah pada sebuah bilangan irasional).
(b) Lengkapilah \cref{ex:b3:topology:sinecurve}: bahwa tak ada \omterm{def:b3:topology:pathconnected}{lintasan} yang menghubungkan
$(0,0)$ dengan $\Gamma$. \emph{(Jika $\gamma = (\gamma_1, \gamma_2)$
sebuah \omterm{def:b3:topology:pathconnected}{lintasan} demikian, misalkan $t^* = \sup\{t : \gamma_1(t) = 0\}$; untuk $t
> t^*$ titiknya bergerak pada grafiknya; pilihlah $t_k \downarrow
t^*$ dengan $\gamma_1(t_k) \to 0$ yang mengenai absis tempat $\sin$
bergantian bernilai $\pm1$ --- teorema nilai antara
menyediakannya --- sehingga bertentangan dengan \omterm{def:b3:topology:continuity}{kekontinuan} $\gamma_2$ di
$t^*$.)}
\end{exercise}

\begin{exercise}[$\star\star\star$]\label{exo:b3:topology:10}
Himpunan Cantor sepertiga tengah $C = \bigcap_n C_n$, dengan $C_0 =
[0,1]$ dan $C_{n+1}$ membuang sepertiga tengah yang terbuka dari setiap
selang $C_n$.
(a) Tunjukkan $C = \bigl\{\sum_{n\geq1} a_n3^{-n} : a_n \in
\{0,2\}\bigr\}$, dan bahwa $C$ \omterm{def:b3:topology:compact}{kompak}, berinterior kosong,
serta tak punya titik terisolasi (yakni \emph{\omterm{prop:b3:galois:perfect}{sempurna}}).
(b) Tunjukkan bahwa $C$ tak \omterm{def:b3:topology:connected}{terhubung} total.
(c) Tunjukkan bahwa $x \mapsto (a_n(x))_n$ merupakan \omterm{def:b3:topology:continuity}{homeomorfisma} $C \to
\{0,2\}^\N$ (dengan \omterm{def:b3:topology:constructions}{topologi hasil kali} pada ruang diskret berisi dua
titik); lalu simpulkan bahwa $C$ tak terbilang, dan bahwa $C \times C
\cong C$.
\end{exercise}

\begin{exercise}[$\star\star\star$]\label{exo:b3:topology:11}
Proyeksi stereografis: dari kutub utara $N = (0, \dots,
0, 1)$ pada $S^n \subseteq \R^{n+1}$, pemetaan $\sigma(x) =
\frac{(x_1, \dots, x_n)}{1 - x_{n+1}}$ merupakan \omterm{def:b3:topology:continuity}{homeomorfisma} $S^n
\setminus \{N\} \to \R^n$ (tuliskan inversnya secara eksplisit). Simpulkan
bahwa $S^n$ \omterm{def:b3:topology:continuity}{homeomorfik} dengan \omterm{def:b3:topology:locallycompact}{kompaktifikasi satu titik}
$\widehat{\R^n}$, dan bahwa membuang \emph{sembarang} titik $S^n$
menyisakan ruang yang \omterm{def:b3:topology:continuity}{homeomorfik} dengan $\R^n$.
\end{exercise}

\begin{exercise}[$\star\star\star$]\label{exo:b3:topology:12}
(\omterm{ex:b3:topology:sinecurve}{Kurva sinus topolog}) Misalkan
\[
T = \bigl\{(x, \sin\tfrac1x) : 0 < x \leq 1\bigr\}
\cup \bigl(\{0\}\times\intcc{-1}1\bigr) \subseteq \R^2 .
\]
(a) Tunjukkan bahwa $T$ \omterm{def:b3:topology:compact}{kompak}, dan bahwa ia merupakan \omterm{def:b3:topology:interior}{penutup}
bagian grafiknya.
(b) Tunjukkan bahwa $T$ \omterm{def:b3:topology:connected}{terhubung} \emph{(sebab grafiknya \omterm{def:b3:topology:connected}{terhubung}
sebagai peta \omterm{def:b3:topology:continuity}{kontinu}; dan \omterm{def:b3:topology:interior}{penutupnya} tetap \omterm{def:b3:topology:connected}{terhubung})}.
(c) Tunjukkan bahwa $T$ \emph{tidak} \omterm{def:b3:topology:pathconnected}{terhubung lintasan}: tak ada \omterm{def:b3:topology:pathconnected}{lintasan}
\omterm{def:b3:topology:continuity}{kontinu} yang menghubungkan $(0,0)$ dengan $(1, \sin1)$. \emph{(Jika $\gamma =
(\gamma_1, \gamma_2)$ sebuah \omterm{def:b3:topology:pathconnected}{lintasan} demikian, misalkan $t^* = \sup\{t :
\gamma_1(t) = 0\}$; tepat setelah $t^*$, $\gamma_1$ mengambil semua
nilai positif yang kecil (menurut teorema nilai antara), sehingga
$\gamma_2$ berayun antara $\pm1$ pada setiap selang
$(t^*, t^* + \delta)$ --- bertentangan dengan \omterm{def:b3:topology:continuity}{kekontinuan} di $t^*$.)}
(d) Simpulkan bahwa keterhubungan \omterm{def:b3:topology:pathconnected}{lintasan} lebih kuat secara sejati
daripada keterhubungan, lalu tunjukkan bahwa contoh semacam itu tak mungkin
terbuka di $\R^n$: sebab himpunan bagian $\R^n$ yang \emph{terbuka} dan \omterm{def:b3:topology:connected}{terhubung} selalu
\omterm{def:b3:topology:pathconnected}{terhubung lintasan} \emph{(himpunan titik yang dapat dihubungkan ke sebuah titik
acuan oleh \omterm{def:b3:topology:pathconnected}{lintasan} bersifat terbuka sekaligus tertutup di dalam daerah asalnya)}.
\end{exercise}

\section{Soal: himpunan Cantor dan sebuah kurva pengisi ruang}

\begin{problem}[{Soal akhir pekan --- kurva Peano memang ada, dan mengapa
ia bukan homeomorfisma}]\label{pb:b3:topology:1}
Pada 1890 Peano mengejutkan analisis dengan sebuah \emph{surjeksi
\omterm{def:b3:topology:continuity}{kontinu}} $[0,1] \to [0,1]^2$: yaitu kurva yang mengisi sebuah persegi. Kita
membangun satu dengan tangan kosong dari himpunan Cantor $C$ pada
\cref{exo:b3:topology:10} (yang hasilnya boleh dipakai bebas),
lalu membuktikan bahwa pemetaan semacam itu tak mungkin injektif: sebab persegi
bukanlah kurva. Sepanjang soal, unsur $C$ ditulis $x =
\sum_{n \geq 1} \frac{2b_n(x)}{3^n}$ dengan angka $b_n(x) \in
\{0, 1\}$.

\medskip
\noindent\textbf{Bagian I --- Membaca angka secara \omterm{def:b3:topology:continuity}{kontinu}.}
\begin{enumerate}
  \item Tunjukkan bahwa jika $x, y \in C$ memenuhi $\abs{x - y} <
        3^{-m}$, maka $b_n(x) = b_n(y)$ untuk setiap $n \leq m$.
        \emph{(Jika angka pertama yang berbeda adalah $b_k$, maka
        $\abs{x - y} \geq \frac{2}{3^k} - \sum_{j>k}\frac2{3^j}
        = \frac1{3^k}$.)}
  \item Simpulkan bahwa setiap fungsi angka $b_n \colon C \to \{0,
        1\}$ bersifat \omterm{def:b3:topology:continuity}{kontinu}, lalu turunkan kembali \omterm{def:b3:topology:continuity}{homeomorfisma} $C
        \cong \{0,1\}^\N$ pada \cref{exo:b3:topology:10}(c).
\end{enumerate}

\medskip
\noindent\textbf{Bagian II --- Sebuah surjeksi \omterm{def:b3:topology:continuity}{kontinu} $C \to
[0,1]^2$.}
\begin{enumerate}[resume]
  \item Definisikan $g \colon C \to [0,1]$ dengan $g(x) = \sum_{n\geq1}
        b_n(x)\,2^{-n}$ (yakni membaca angka Cantor sebagai biner).
        Tunjukkan bahwa $g$ \omterm{def:b3:topology:continuity}{kontinu} (pakailah pertanyaan 1) dan
        surjektif. Apakah ia injektif?
  \item Definisikan $\Phi \colon C \to [0,1]^2$ dengan
        \[
        \Phi(x) = \Bigl(\ \sum_{n\geq1} b_{2n-1}(x)\,2^{-n},\
        \sum_{n\geq1} b_{2n}(x)\,2^{-n}\ \Bigr)
        \]
        (angka ganjil memberikan absisnya, angka genap
        ordinatnya). Tunjukkan bahwa $\Phi$ \omterm{def:b3:topology:continuity}{kontinu} dan
        surjektif.
\end{enumerate}

\medskip
\noindent\textbf{Bagian III --- Menutup celahnya: kurva
Peano.}
\begin{enumerate}[resume]
  \item Komplemen $[0,1] \setminus C$ merupakan gabungan terbilang
        selang terbuka $(u, v)$ yang saling lepas (yaitu sepertiga tengah
        yang dibuang) dan ujungnya terletak di $C$. Definisikan $f \colon
        [0,1] \to [0,1]^2$ dengan $f = \Phi$ pada $C$, lalu diperluas
        secara \emph{afin} pada setiap celahnya:
        \[
        f(t) = \frac{v - t}{v - u}\,\Phi(u) +
        \frac{t - u}{v - u}\,\Phi(v),
        \qquad t \in (u, v).
        \]
        Tunjukkan bahwa $f$ terdefinisi dengan baik dan surjektif pada
        $[0,1]^2$.
  \item Tunjukkan bahwa $f$ \omterm{def:b3:topology:continuity}{kontinu} di setiap titik $[0,1]
        \setminus C$ (sebab afin secara lokal), dan di setiap titik
        $C$: diberikan $\varepsilon$, pilihlah $m$ dengan $2^{-m} <
        \varepsilon/2$ lalu pakai pertanyaan 1 untuk mengendalikan $\Phi$ pada
        $C$ di dekat $x$, lalu periksalah bahwa interpolasi
        afinnya tak dapat lolos: pada sebuah celah $(u,v) \subseteq
        (x - 3^{-m}, x + 3^{-m})$, nilai $f(t)$ terletak pada
        ruas garis $[\Phi(u), \Phi(v)]$, yang kedua ujungnya dekat dengan
        $\Phi(x)$. Simpulkan: $f$ merupakan \emph{surjeksi
        \omterm{def:b3:topology:continuity}{kontinu}} $[0,1] \to [0,1]^2$.
  \item Simpulkan adanya surjeksi \omterm{def:b3:topology:continuity}{kontinu} $[0,1] \to [0,1]^n$ untuk
        setiap $n \geq 2$, dan juga $\R \to \R^2$.
\end{enumerate}

\medskip
\noindent\textbf{Bagian IV --- Namun tak pernah injektif.}
\begin{enumerate}[resume]
  \item Tunjukkan bahwa \emph{injeksi} \omterm{def:b3:topology:continuity}{kontinu} $f \colon [0,1]
        \to [0,1]^2$ akan menjadi \omterm{def:b3:topology:continuity}{homeomorfisma} pada petanya
        (\cref{cor:b3:topology:compacthomeo}).
  \item Tunjukkan bahwa $[0,1]$ dan $[0,1]^2$ tidak \omterm{def:b3:topology:continuity}{homeomorfik}:
        buanglah sebuah titik yang dipilih baik-baik lalu bandingkan keterhubungannya
        (\cref{met:b3:topology:connectargument}).
  \item Simpulkan: surjeksi \omterm{def:b3:topology:continuity}{kontinu} $[0,1] \to [0,1]^2$
        tak pernah dapat injektif --- sebab yang injektif akan membuat
        $[0,1]^2 = f([0,1])$ \omterm{def:b3:topology:continuity}{homeomorfik} dengan $[0,1]$ menurut
        pertanyaan 8, dan itu bertentangan dengan pertanyaan 9. Di manakah tepatnya
        kekompakan $[0,1]$ masuk ke dalam argumennya?
  \item (Puncaknya) Rangkailah pelajarannya: ada surjeksi
        \omterm{def:b3:topology:continuity}{kontinu} tetapi tak ada bijeksi \omterm{def:b3:topology:continuity}{kontinu} $[0,1] \to
        [0,1]^2$. Apa yang dikatakan hal ini tentang ``dimensi'' sebagai gagasan
        \omterm{def:b3:topology:topology}{topologi}? Rumuskan secara tepat satu teorema yang
        dibuktikan pada soal ini dan satu pernyataan masuk akal yang
        tetap di luar jangkauan perkakas kita (yakni invariansi domain).
\end{enumerate}

\medskip
\noindent\textbf{Bagian V --- Aritmetika himpunan Cantor.}
\begin{enumerate}[resume]
  \item Buktikan bahwa $C + C = [0, 2]$: diberikan $t \in [0, 1]$,
        tulis $t = \sum_n d_n3^{-n}$ dengan angka $d_n \in \{0,
        1, 2\}$ lalu pecahlah setiap angkanya sebagai $d_n = a_n + b_n$ dengan
        $a_n, b_n \in \{0, 1\}$; simpulkan bahwa $\frac C2 +
        \frac C2 = [0, 1]$ \emph{(himpunan bagian $[0,1]$ manakah
        $\frac C2$ itu, dinyatakan dalam angka ternernya?)}, lalu skalakan kembali.
        Tak pernah ada simpanan yang diperlukan --- katakan mengapa itulah intinya.
  \item Tafsirkan secara geometri: ``debu Cantor'' $C \times
        C \subseteq \R^2$, yang berinterior kosong, melemparkan
        bayangan \emph{penuh} pada diagonalnya: sebab proyeksi $(x,
        y) \mapsto x + y$ memetakannya pada $[0, 2]$. Catat pula
        keserupaan dirinya $C = \frac C3 \cup \bigl(\frac C3 +
        \frac23\bigr)$, yaitu persamaan di balik setiap gambar
        $C$.
  \item Tunjukkan bahwa penyelipan angka berselang mendefinisikan \omterm{def:b3:topology:continuity}{homeomorfisma} $C
        \to C \times C$, lalu simpulkan $C \cong C^n$ untuk setiap $n$
        --- jadi himpunan Cantor adalah kuadrat dirinya sendiri, juga kubusnya, \dots{}
        Ruang akrab manakah yang berbagi sifat ini?
  \item Tunjukkan bahwa $\frac14 \in C$ \emph{(hitunglah ekspansi
        ternernya: $\frac14 = \sum_{k\geq1} 2\cdot3^{-2k}$)}
        walaupun $\frac14$ bukan ujung selang mana pun yang
        dibuang; lalu simpulkan --- dengan mencacah ujungnya --- bahwa
        ujung-ujung itu membentuk himpunan bagian sejati $C$ yang terbilang.
  \item Tunjukkan bahwa pemetaan pembacaan biner $g$ pada pertanyaan 3 bersifat
        paling banyak $2$-ke-$1$, lalu uraikan dengan tepat titik mana
        di $[0,1]$ yang mempunyai dua prapeta. (Pemetaan $g$ meruntuhkan $C$ pada
        $[0,1]$ dengan merekatkan terbilang banyak pasangan: yaitu bayangan
        kombinatorik tangga Cantor yang akan dijumpai lagi
        di \cref{ch:b3:measure}.)
\end{enumerate}

\medskip
\noindent\textbf{Bagian VI --- Himpunan \omterm{def:b3:topology:compact}{kompak} \omterm{prop:b3:galois:perfect}{sempurna}: Cantor
di mana-mana.} Sebuah ruang metrik \omterm{def:b3:topology:compact}{kompak} tak kosong $P$ disebut
\emph{\omterm{prop:b3:galois:perfect}{sempurna}} jika ia tak punya titik terisolasi.
\begin{enumerate}[resume]
  \item Misalkan $P$ \omterm{def:b3:topology:compact}{kompak} \omterm{prop:b3:galois:perfect}{sempurna} dan misalkan $x \in P$, $r > 0$.
        Tunjukkan bahwa bola $B(x, r)$ memuat dua titik $y_0
        \neq y_1$ dari $P$, dan karenanya dua bola tertutup
        \emph{saling lepas} $\bar B(y_0, \rho)$, $\bar B(y_1, \rho)$ di dalam
        $B(x, r)$, yang berpusat di titik $P$ dan berjari-jari $\rho
        \leq r/4$ sekecil yang diinginkan. Jelaskan mengapa kesempurnaannya
        (yakni tanpa titik terisolasi) tepat merupakan syarat yang membuat pemecahan
        ini dapat diulang di dalam \emph{masing-masing} bola
        baru itu.
  \item Iterasikan: bangunlah himpunan tertutup $F_w \neq \varnothing$
        yang diindeks kata biner berhingga $w$, dengan $F_{w0},
        F_{w1} \subseteq F_w$ saling lepas dan
        $\operatorname{diam}F_w \leq 2^{-\abs w}
        \operatorname{diam}P$. Tunjukkan bahwa untuk setiap kata tak
        berhingga $\varepsilon \in \{0,1\}^\N$ irisan
        $\bigcap_mF_{\varepsilon_1\cdots\varepsilon_m}$ berupa
        satu titik $h(\varepsilon)$ \emph{(lewat sifat irisan
        berhingga pada $P$ yang \omterm{def:b3:topology:compact}{kompak})}.
  \item Tunjukkan bahwa $h\colon\{0,1\}^\N \to P$ injektif dan
        \omterm{def:b3:topology:continuity}{kontinu}, lalu simpulkan: \emph{setiap ruang metrik \omterm{def:b3:topology:compact}{kompak}
        yang \omterm{prop:b3:galois:perfect}{sempurna} bersifat tak terbilang} --- bahkan berkardinalitas
        sekurang-kurangnya sebesar $\{0,1\}^\N$. Peroleh kembali: $[0,1]$ dan $C$
        tak terbilang.
  \item Tunjukkan bahwa $h$ merupakan \omterm{def:b3:topology:continuity}{homeomorfisma} pada petanya
        \emph{(sebab injeksi \omterm{def:b3:topology:continuity}{kontinu} dari sebuah \omterm{def:b3:topology:compact}{kompak},
        \cref{cor:b3:topology:compacthomeo})}: jadi setiap ruang metrik
        \omterm{def:b3:topology:compact}{kompak} yang \omterm{prop:b3:galois:perfect}{sempurna} memuat salinan \omterm{def:b3:topology:continuity}{homeomorfik} dari
        himpunan Cantor. Himpunan Cantor bukanlah keanehan melainkan
        benih universal bagi kesempurnaan \omterm{def:b3:topology:compact}{kompak}.
  \item Simpulkan bahwa setiap ruang metrik \omterm{def:b3:topology:compact}{kompak} yang terbilang mempunyai
        titik terisolasi, lalu tampilkan satu contoh yang titik
        terisolasinya \omterm{def:b3:topology:interior}{padat} tetapi bukan seluruhnya: $K = \{0\} \cup
        \{\frac1n : n \geq 1\}$.
  \item (Cantor--Bendixson untuk $\R$) Misalkan $F \subseteq \R$
        tertutup. Sebutlah $x$ sebuah \emph{titik kondensasi} $F$ jika
        setiap \omterm{def:b3:topology:topology}{persekitaran} $x$ memotong $F$ secara tak terbilang. Tunjukkan
        bahwa titik kondensasi sebuah $F$ tertutup yang tak terbilang
        membentuk himpunan tertutup \omterm{prop:b3:galois:perfect}{sempurna} tak kosong $P$, dan bahwa $F
        \setminus P$ terbilang \emph{(selimuti titik
        yang bukan kondensasi dengan terbilang banyak selang rasional yang memotong
        $F$ secara terbilang)}. Simpulkan: setiap
        himpunan bagian $\R$ yang tertutup bersifat terbilang atau berkardinalitas
        kontinum --- jadi hipotesis kontinum berlaku
        untuk himpunan tertutup.
\end{enumerate}

\medskip
\noindent\textbf{Bagian VII --- \omterm{def:b3:topology:interior}{Penutup}: seberapa mulus, dan seberapa
jauh dari \omterm{prop:b3:galois:perfect}{sempurna}.}
\begin{enumerate}[resume]
  \item (Kemulusan Hölder pada kurvanya) Tetapkan $\alpha =
        \frac{\ln 2}{2\ln 3} \approx 0.3155$, lalu catat
        identitas $3^\alpha = \sqrt2$. Dengan memakai pertanyaan 1, tunjukkan
        bahwa
        \[
        \norm{\Phi(x) - \Phi(y)}_\infty \leq 2\,\abs{x -
        y}^\alpha \qquad (x, y \in C),
        \]
        lalu rambatkan taksirannya melalui celah afinnya:
        tunjukkan bahwa kurva Peano $f$ pada pertanyaan 6 memenuhi
        $\norm{f(t) - f(s)}_\infty \leq 6\abs{t - s}^\alpha$ pada
        seluruh $[0,1]$ \emph{(tangani pasangan di celah yang sama lewat
        interpolasi, lalu pasangan umumnya lewat titik ekstrem $C \cap [t, s]$)}.
  \item (Eksponen $\frac12$ adalah tembok) Tunjukkan bahwa tak ada
        surjeksi $F \colon [0,1] \to [0,1]^2$ yang dapat bersifat
        Hölder-$\alpha$ dengan $\alpha > \frac12$: potonglah $[0,1]$
        menjadi $N$ selang, batasilah diameter
        petanya, lalu cacahlah titik kisi
        $\bigl(\frac iM, \frac jM\bigr)$, $0 \leq i, j \leq M$,
        yang dapat dimuat sebuah himpunan berdiameter $< \frac1M$; lalu pilihlah
        $M$ berorde $N^\alpha$ dan biarkan $N \to \infty$. Letakkan
        kurva kita ($\alpha \approx 0.3155$) relatif terhadap tembok itu, lalu
        catat tanpa bukti bahwa kurva Hilbert
        mencapai eksponen kritis $\frac12$.
  \item (Himpunan turunan: mengukur ketaksempurnaan) Untuk $F$ yang tertutup
        di $\R$, misalkan $F'$ himpunan titik limit $F$ (yang juga
        tertutup), lalu iterasikan: $F^{(0)} = F$, $F^{(i+1)}
        = (F^{(i)})'$. Periksalah bahwa
        \[
        K = \{0\} \cup \Bigl\{\frac1m\Bigr\}_{m \geq 1} \cup
        \Bigl\{\frac1m + \frac1n : m \geq 1,\ n > m^2\Bigr\}
        \]
        merupakan himpunan \omterm{def:b3:topology:compact}{kompak} terbilang dengan $K' = \{0\} \cup
        \{\frac1m\}$, $K'' = \{0\}$, $K''' = \varnothing$
        \emph{(periksalah bahwa gugusan ke-$m$ hidup di dalam
        selang $(\frac1m, \frac1{m-1})$, sehingga gugusannya tidak
        saling menyisip)}, dan bahwa titik terisolasinya \omterm{def:b3:topology:interior}{padat}
        di $K$, seperti diramalkan pertanyaan 21. Uraikan
        induksi yang menghasilkan, untuk setiap $j$, sebuah \omterm{def:b3:topology:compact}{kompak}
        terbilang $K_j$ dengan $K_j^{(j)} = \{0\}$ dan $K_j^{(j+1)} =
        \varnothing$, lalu bandingkan dengan himpunan \omterm{prop:b3:galois:perfect}{sempurna} pada
        pertanyaan 22, yang penurunannya tak pernah bergerak: jadi
        rank berhingga persis merupakan lawan kesempurnaan.

\end{enumerate}
\end{problem}
